% ===================================================================
%  तक्ता क्रमांक ३ — बालपण आणि तारुण्य.
%
%  Traduction, faite en 2026, en marathi d'aujourd'hui, sur l'ido de
%  texto/io. Regles de la colonne : voir 10-takta-01.tex. Deux
%  scenes, deux numerotations : les cles portent c1 et c2.
%
%  LE MARATHI POSTPOSE SES RAPPORTS, ET C'EST LA MEME CONTRAINTE QUE
%  LE CHINOIS. « घंटाघराचा कळस » — du clocher, la fleche — met le
%  possesseur AVANT le possede, la ou l'ido fait l'inverse ; chaque
%  « X de Y » sort donc les renvois a l'envers. Le remede est celui
%  qu'on avait mis au point pour le cantonais : nommer d'abord ce que
%  l'ido nomme d'abord, puis qualifier par une seconde proposition.
%  Le bloc c1-04-1 — la fleche du clocher de l'eglise — en est le
%  premier cas de la colonne, et il a ete ecrit ainsi du premier jet.
%
%  Ce n'est pas une nouveaute pour la famille : renvoji.py porte dans
%  son en-tete l'exemple hindi « अध्यापक (7) अपने मंच (6) पर », et le
%  hindi avait coute onze inversions a lui seul. L'arabe et l'ido
%  antposent, le chinois et le marathi postposent — deux camps, et la
%  colonne change de camp en changeant de langue.
%
%  LA FETE DE VILLAGE PORTE SES NOMS MARATHIS, ET DEUX TOMBENT JUSTE :
%
%    जत्रा (39)   la foire — la fete patronale du village, avec ses
%                  manèges, son cirque et sa ménagerie
%    सरपंच (37)  le maire — le chef elu d'un village
%
%  Ce sont les memes rencontres que le مولد et la عمدة egyptiens au
%  meme tableau : un seul mot pour la chose entiere, sans periphrase.
%
%  ET LE RESTE DU VILLAGE :
%
%    साखरफुटाणे (55) les dragees qu'on jette aux enfants
%    खडावा (75)      les sabots de bois
%    पाकोळी (4)      l'hirondelle
%    मोहळ (21)       la ruche
%    फुलपाखरू (35)   le papillon  हिन्दी : तितली
%    घसरगुंडी (49)   le toboggan
%    आखाडा (44)      l'arene
%    ढग (2)          le nuage     हिन्दी : बादल
%    गाढव (23)       l'ane        हिन्दी : गधा
%    म्हातारी (71)   la vieille femme
%    भिकारी (66)     le mendiant
%    शेतकरी (69)     le paysan
%
%  ET L'HABILLEMENT EST TOUT MARATHI : चोळी (29) le corsage, परकर
%  (31) le jupon, बंडी (17) le gilet, विजार (16) le pantalon, सदरा
%  (12) la chemise, मोजे (6) les chaussettes, चड्डी (3) le calecon,
%  बनियन (2) le maillot de corps.
% ===================================================================

%%K t03-apar-1 apar
\VUsaut{3.0mm}
\VUtitre{15.0pt}{60}{तक्ता क्रमांक ३}
\VUsaut{1.6mm}
\VUtitre{11.4pt}{0}{बालपण आणि तारुण्य.}
\VUsaut{3.4mm}

%%K t03-apar-2 apar
(तक्ता क्रमांक ३ आणि ४ असे जोडलेले आहेत की ते चार \VUgras{कौटुंबिक प्रवेशांतून}
\VUgras{हवा}, \VUgras{वय}, \VUgras{कुटुंब}, \VUgras{पोशाख} आणि \VUgras{स्थिती}
यांच्या मुख्य कल्पना मांडतात.)

\VUsaut{4.0mm}
%%K t03-tit-1 sub
\VUcentre{12.0pt}{0}{\textit{पहिला प्रवेश.}}
\VUsaut{3.0mm}

%%K t03-tit-2 sub
\VUpk{10.2pt}{40}{गावातला बाप्तिस्मा-सोहळा.}
\VUsaut{2.4mm}

%%K t03-tit-3 sub
\VUpk{10.2pt}{40}{वसंत ऋतू.}
\VUsaut{3.0mm}

%%K t03-c1-01-1 p
१. --- आपण \VUgras{वसंत ऋतूत} आहोत, \VUgras{मार्च}, \VUgras{एप्रिल} किंवा
\VUgras{मे} महिन्यात, \VUgras{आठवड्यातल्या} \VUgras{सोमवार},
\VUgras{मंगळवार} किंवा \VUgras{बुधवार} या दिवशी. सकाळचे \VUgras{दहा वाजले}
आहेत.

%%K t03-c1-02-1 p
२. --- \VUgras{हवा} सुंदर आहे, \VUgras{वारा} स्वच्छ आहे. सूर्य तळपतो आहे.
काही लहान \VUgras{ढग} \textsuperscript{(2)} वर चढत आहेत --- निळ्या
\VUgras{आकाशात} \textsuperscript{(1)}.

%%K t03-c1-03-1 p
३. --- \VUgras{झाडांना} अजून जेमतेम \VUgras{पाने} आहेत. सडपातळ
\VUgras{चिनार} \textsuperscript{(3)} आणि उंच \VUgras{प्लेन वृक्ष}
\textsuperscript{(17)} यांना आतापर्यंत फक्त कळ्याच आहेत.

%%K t03-c1-04-1 p
४. --- \VUgras{पाकोळ्या} \textsuperscript{(4)} परत आल्या आहेत आणि आनंदाने
किलबिलत \VUgras{कळसाभोवती} \textsuperscript{(7)} घिरट्या घालत आहेत --- तो कळस
\VUgras{घंटाघराचा} \textsuperscript{(5)} आहे, आणि ते घंटाघर गावच्या
\VUgras{चर्चवर} \textsuperscript{(6)} मुकुटासारखे उभे आहे.

%%K t03-tit-4 sub
\VUsaut{3.4mm}
\VUpk{10.2pt}{40}{चर्च.}
\VUsaut{3.0mm}

%%K t03-c1-05-1 p
५. --- हे चर्च फार साधे आहे, त्याच्या मोठ्या \VUgras{महाद्वारासह}
\textsuperscript{(13)} आणि जाड \VUgras{घड्याळासह} \textsuperscript{(8)}.
\VUgras{लहान काटा} \textsuperscript{(9)} X या आकड्यावर आहे; तो \VUgras{तास}
दाखवतो.

%%K t03-c1-05-1 p suite
\VUgras{मोठा} \textsuperscript{(10)} XII वर आहे (मध्यान्ह); तो \VUgras{मिनिटे}
दाखवतो.

%%K t03-c1-06-1 p
६. --- घंटाघरात \VUgras{घंटा} \textsuperscript{(11)} आनंदाने वाजत आहेत. छपराच्या
शिखरावर दगडाचा एक लहान \VUgras{क्रूस} \textsuperscript{(12)} उभा आहे.

%%K t03-c1-07-1 p
७. --- चर्चला फक्त मोठे \VUgras{महाद्वारच} \textsuperscript{(13)} नाही, बाजूला एक
लहान दारही आहे. या लहान दारातून \VUgras{धर्मगुरू} \textsuperscript{(15)} आपला
\VUgras{झगा} घालून \VUgras{वस्त्रगृहातून} \textsuperscript{(14)} बाहेर पडतात आणि
\VUgras{धर्मगुरूच्या घराकडे} \textsuperscript{(19)} जातात.

%%K t03-c1-08-1 p
८. --- त्यांच्याजवळ, अहो, \VUgras{दूधवाली} \textsuperscript{(24)} आपल्या
\VUgras{गाढवासह} \textsuperscript{(23)} आहे. ती शहरातून परत येते आहे, जिथे तिने
मुलांसाठी चांगले दूध नेले होते.

%%K t03-tit-5 sub
\VUsaut{4.0mm}
\VUpk{10.2pt}{40}{फळबाग.}
\VUsaut{3.4mm}

%%K t03-c1-09-1 p
९. --- अलिबोरॉनराव (गाढव) या सकाळच्या फेरीने अगदी खूश आहेत. शेजारच्या
\VUgras{फळबागेतल्या} \VUgras{फुलांच्या} \textsuperscript{(20)} सुगंधाने भरलेली हवा
ते मजेत हुंगतात --- ती फुले \VUgras{फळझाडांना} \textsuperscript{(18)} झाकून
टाकतात. ती \VUgras{पीचची झाडे} \textsuperscript{(18)} आणि \VUgras{जर्दाळूची झाडे}
\textsuperscript{(19)} अगदी गुलाबी आणि पांढरी झाली आहेत, आणि चांगल्या पिकाची आशा
देतात.

%%K t03-c1-10-1 p
१०. --- त्याच वेळी लहान \VUgras{मधमाश्या} \textsuperscript{(22)} --- त्या
\VUgras{मोहळातल्या} \textsuperscript{(21)} आहेत --- गुणगुणत तिथे आपला गोड
\VUgras{मध} गोळा करायला जातात.

%%K t03-c1-11-1 p
११. --- पण हे काय? अहो, गावातली मुले धावत येत आहेत आणि घाबरलेल्या \VUgras{हंसांना}
\textsuperscript{(25)} पळवून लावत आहेत. वकिलाच्या मुलीच्या \VUgras{बाप्तिस्म्यानंतर}
मंडळी चर्चमधून बाहेर पडते आहे.

%%K t03-c1-12-1 p
१२. --- लहानगा जेम्स आपल्या \VUgras{पाळण्यात} \textsuperscript{(26)} घंटांच्या
आनंदी नादाने जागा होतो, आणि त्याची बहीण मॅग्डलीन --- तिलाही बाप्तिस्म्याची
मिरवणूक पाहायची आहे --- आपल्या आईला विनवते की तिने घराच्या उंबऱ्यावर तिचे केस
विंचरावेत आणि तिला कपडे घालावेत.

%%K t03-c1-13-1 p
१३. --- आई होकार देते. ती आपली \VUgras{ओढणी} \textsuperscript{(28)}, आपली
\VUgras{चोळी} \textsuperscript{(29)} आणि आपला \VUgras{एप्रन} \textsuperscript{(30)}
घालते, आणि दारासमोर एका लहान खुर्चीवर येऊन बसते. मॅग्डलीनच्या अंगावर अजून फक्त
तिचा \VUgras{परकर} \textsuperscript{(31)} आणि तिची \VUgras{कंचुकी}
\textsuperscript{(32)} आहे.

%%K t03-c1-14-1 p
१४. --- ती किती आनंदी आहे! तिला आपली आवडती फुले विसरून जातात --- तिची
\VUgras{कार्नेशने} \textsuperscript{(34)}, ज्यांवर \VUgras{फुलपाखरे}
\textsuperscript{(35)} बसतात, तिचे \VUgras{ट्यूलिप} \textsuperscript{(36)}, तिची
\VUgras{खोऱ्यातली कमळे} \textsuperscript{(37)}, जी \VUgras{कुंड्यांत}
\textsuperscript{(38)} \VUgras{भिंतीवर} \textsuperscript{(33)} आहेत, आणि तिचे
\VUgras{गुलाब} \textsuperscript{(39)}, ज्यांचे इतके सुंदर कुंपण झाले आहे. ती
मिरवणुकीतल्या बायांचे श्रीमंत कपडे न्याहाळते आहे.

%%K t03-c1-15-1 p
१५. --- गावातली मुले कपड्यांकडे फारसे लक्ष देत नाहीत. \VUgras{साखरफुटाणे}
\textsuperscript{(55)} किंवा \VUgras{नाणी} \textsuperscript{(52)} मिळावीत म्हणून ती
गर्दी करतात आणि एकमेकांना धक्के देतात. त्यांतले एक रडते आहे
\textsuperscript{(40)}.

%%K t03-c1-16-1 p
१६. --- दुसऱ्याने \VUgras{कुत्र्याच्या} \textsuperscript{(42)} पायावर पाय दिला;
तो ओरडत पळाला आणि त्याने \VUgras{कोंबडीला} \textsuperscript{(43)} व तिच्या
\VUgras{पिल्लांना} \textsuperscript{(44)} पळवून लावले.

%%K t03-tit-6 sub
\VUsaut{5.0mm}
\VUpk{10.2pt}{40}{बाप्तिस्म्याची मिरवणूक.}
\VUsaut{4.0mm}

%%K t03-c1-17-1 p
१७. --- मिरवणुकीच्या पुढे \VUgras{दाई} \textsuperscript{(45)} चालते आहे. तिच्या
डोक्यावर लांब फिती असलेली सुंदर \VUgras{टोपी} \textsuperscript{(46)} आहे. ती
\VUgras{नवजात बाळाला} \textsuperscript{(47)} घेऊन चालली आहे, ज्याच्या अंगावर
सगळीकडे \VUgras{लेस} \textsuperscript{(48)} आहे.

%%K t03-c1-18-1 p
१८. --- दाईच्या मागे \VUgras{धर्मपिता} \textsuperscript{(49)} आणि
\VUgras{धर्ममाता} \textsuperscript{(53)} येतात. तेच साखरफुटाणे आणि नाणी वाटतात.
धर्मपित्याच्या हातात \VUgras{हातमोजे} \textsuperscript{(51)} आहेत आणि डोक्यावर
\VUgras{उंच टोपी} \textsuperscript{(50)} आहे. धर्ममातेच्या हातात सुंदर
\VUgras{फुलांचा गुच्छ} \textsuperscript{(54)} आहे.

%%K t03-c1-19-1 p
१९. --- त्यांच्या मागे आपल्याला बाळाचे \VUgras{काका} \textsuperscript{(56)} आणि
\VUgras{काकू} \textsuperscript{(58)} दिसतात. काकूंच्या डोक्यावर सुरेख
\VUgras{टोपी} \textsuperscript{(59)} आहे, आणि काकांच्या अंगावर काळा \VUgras{लांब
कोट} \textsuperscript{(57)} व पांढरे जाकीट आहे. मग \VUgras{चुलतभाऊ}
\textsuperscript{(60)} आणि \VUgras{चुलतबहीण} \textsuperscript{(61)}. ही
नवजाताच्या \VUgras{बहिणीचा} \textsuperscript{(62)} हात धरून आहे --- लहानगी बर्था.

%%K t03-c1-20-1 p
२०. --- बर्थाचा दुसरा \VUgras{भाऊ} \textsuperscript{(63)} मागून चालतो आहे. तो एका
\VUgras{नातेवाईक बाईचा} \textsuperscript{(64)} हात धरून आहे. कुटुंबाचे
\VUgras{मित्र} \textsuperscript{(65)} त्यानंतर येतात.

%%K t03-tit-7 sub
\VUsaut{4.6mm}
\VUpk{10.2pt}{40}{बघे.}
\VUsaut{3.4mm}

%%K t03-c1-21-1 p
२१. --- मिरवणुकीजवळ, अहो, एक \VUgras{भिकारी} \textsuperscript{(66)} आपल्या
\VUgras{कुबडीवर} \textsuperscript{(67)} टेकून उभा आहे. तो भीक मागतो आहे.

%%K t03-c1-22-1 p
२२. --- दुसऱ्या बाजूला एक फार \VUgras{सभ्य} \textsuperscript{(68)} लहान मुलगा आहे.
तो आपल्या ओळखीच्या माणसांना नमस्कार करतो आणि «\VUgras{सुप्रभात}» म्हणतो.

%%K t03-c1-23-1 p
२३. --- बाप्तिस्म्याची मिरवणूक पाहायला गावकरी आपापल्या घरांतून बाहेर पडले आहेत.
आपल्याला एक \VUgras{शेतकरी} \textsuperscript{(69)} त्याच्या \VUgras{सुती टोपीसह}
दिसतो, आणि त्याच्या शेजारी एक \VUgras{शेतकरीण} \textsuperscript{(70)}. मग एक
\VUgras{म्हातारी} \textsuperscript{(71)} आपल्या \VUgras{काठीवर}
\textsuperscript{(72)} टेकलेली. शेवटी \VUgras{गिरणीवाला} \textsuperscript{(73)}
आपल्या मोठ्या \VUgras{लोकरी टोपीसह} \textsuperscript{(74)}, आणि एक तरुण शेतकरीण
\VUgras{खडावा} \textsuperscript{(75)} घालून, जी आपल्या तान्ह्या मुलाला चालायला
शिकवते आहे.

%%K t03-c1-24-1 p
२४. --- हा प्रवेश फार मजेदार आहे.

\VUsaut{4.0mm}
%%K t03-tit-8 sub
\VUcentre{12.0pt}{0}{\textit{दुसरा प्रवेश.}}
\VUsaut{3.0mm}

%%K t03-tit-9 sub
\VUpk{10.2pt}{40}{सार्वजनिक उत्सव.}
\VUsaut{2.4mm}

%%K t03-tit-10 sub
\VUpk{10.2pt}{40}{नौकाक्रीडा आणि शर्यती.}
\VUsaut{3.0mm}

%%K t03-c2-01-1 p
१. --- \VUgras{उन्हाळा} आला आहे. \VUgras{जून} (जुलै किंवा \VUgras{ऑगस्ट})
महिन्याचा तापता सूर्य तरुणांना पोहायला आणि नाव वल्हवायला उद्युक्त करतो.

%%K t03-c2-02-1 p
२. --- नदीच्या उजव्या तीरावर नावाडी नौकाक्रीडेत आणि शर्यतींत भाग घेण्यासाठी
आपले कपडे काढत आहेत.

%%K t03-c2-03-1 p
३. --- त्यांतला एक आपले \VUgras{बूट} \textsuperscript{(1)} काढतो आहे; तो त्यांच्या
नाड्या सोडतो आहे. त्याचे दोन सोबती तयार झाले आहेत. आता त्यांच्या अंगावर फक्त
त्यांचे \VUgras{बनियन} \textsuperscript{(2)} आणि \VUgras{पोहण्याची चड्डी}
\textsuperscript{(3)} आहे. आपल्या मित्रांची वाट पाहत ते गप्पा मारत आहेत. ते
जत्रेविषयी आणि पलीकडच्या तीरावर चाललेल्या उत्सवाविषयी बोलत आहेत.

%%K t03-c2-04-1 p
४. --- त्यांच्याजवळ जमिनीवर त्यांच्या सोबत्यांचे \VUgras{कपडे} पडले आहेत: एक जोड
\VUgras{बुटांचा} \textsuperscript{(4)}, \VUgras{वहाणा} \textsuperscript{(5)},
\VUgras{मोजे} \textsuperscript{(6)}, \VUgras{लोकरी} \textsuperscript{(7)} आणि
\VUgras{गवती} \textsuperscript{(8)} टोप्या, आणि एक \VUgras{नेकटाय}
\textsuperscript{(9)}.

%%K t03-c2-05-1 p
५. --- तरुणांतल्या एकाने आपला \VUgras{सूट} \textsuperscript{(10)} काळजीपूर्वक घडी
घालून ठेवला आहे. तो त्याला एका टॉवेलवर ठेवतो आहे, ज्यात त्याचा शेजारी लवकरच दोन्ही
सूट गुंडाळणार आहे.

%%K t03-c2-06-1 p
६. --- सहावा मुलगा उशीर करतो आहे. त्याच्या डोक्यावर अजून त्याची \VUgras{कॅप}
\textsuperscript{(11)} आहे. त्याने आपला \VUgras{सदरा} \textsuperscript{(12)}, आपली
\VUgras{कॉलर} \textsuperscript{(13)} किंवा आपले \VUgras{कफ} \textsuperscript{(14)}
काढलेले नाहीत. हातात तो आपली \VUgras{बंडी} \textsuperscript{(17)} धरून आहे, आणि
अंगावर अजून त्याची \VUgras{विजार} \textsuperscript{(16)} आणि त्याचे \VUgras{गॅलिस}
\textsuperscript{(15)} आहेत. पुढच्या शर्यतीसाठी तो नक्कीच तयार होणार नाही.

%%K t03-c2-07-1 p
७. --- तरुणांजवळून एक पाऊलवाट जाते, जिच्या दोन्ही बाजूंना पुष्कळ \VUgras{काटेरी
झुडपे} \textsuperscript{(18)} आहेत, आणि ती नदीच्या काठापर्यंत नेते, जिथे जेकब
काका \VUgras{लव्हाळ्यांमध्ये} \textsuperscript{(19)} \VUgras{शोभेची दारू}
\textsuperscript{(20)} तयार करत आहेत --- ती संध्याकाळी उडवली जाईल.

%%K t03-tit-11 sub
\VUsaut{4.4mm}
\VUpk{10.2pt}{40}{शर्यत.}
\VUsaut{3.4mm}

%%K t03-c2-08-1 p
८. --- नदीवर काही बाया आपल्या छत्र्यांखाली सावली धरून \VUgras{नावेतून}
\textsuperscript{(21)} फेरफटका मारत आहेत. त्या शर्यत पाहत आहेत, जी फार रोचक आहे.
प्रत्येक \VUgras{होडीत} \textsuperscript{(22)} चार \VUgras{वल्हवणारे}
\textsuperscript{(24)} आपल्या सगळ्या ताकदीने वल्हवत आहेत,

%%K t03-c2-08-1 p suite
आणि त्याच वेळी \VUgras{सुकाणू धरणारा} \textsuperscript{(26)} \VUgras{सुकाणू}
\textsuperscript{(25)} धरून नाजूक होडी वळवतो आहे. \VUgras{वल्ही}
\textsuperscript{(23)} तालावर पाणी कापत आहेत आणि हलक्या होड्या चक्रावून टाकणाऱ्या
वेगाने पळत आहेत.

%%K t03-c2-09-1 p
९. --- पुलाच्या खांबाजवळ काही तरुण \VUgras{निसरड्या खांबाच्या}
\textsuperscript{(28)} बक्षिसासाठी झुंजत आहेत. त्यांतला एक लवचिक व निसरड्या
खांबावरून सावधपणे चालतो आहे --- टोकाला बांधलेल्या लहान \VUgras{झेंड्यापर्यंत}
\textsuperscript{(29)} पोहोचण्यासाठी. त्याचा दुर्दैवी \VUgras{स्पर्धक}
\textsuperscript{(30)} पाण्यात पडला आहे. किनाऱ्यावर येण्यासाठी तो पोहतो आहे.

%%K t03-c2-10-1 p
१०. --- अधिक वयाचे तरुण \VUgras{नौका-लढत} \textsuperscript{(32)} खेळत आहेत ---
\VUgras{धक्क्याजवळ} \textsuperscript{(33)}. या सगळ्या खेळांना मोठी \VUgras{गर्दी}
\textsuperscript{(31)} पाहत आहे --- \VUgras{पुलाच्या} \textsuperscript{(27)}
\VUgras{कठड्याला} टेकलेली आणि \VUgras{तीरावर} जमलेली. तरीही काही बघे
\VUgras{बक्षिसाच्या खांबाचा} \textsuperscript{(45)} खेळ पाहायला वळतात, किंवा
बक्षिसे \VUgras{वाटायला} येणारी \VUgras{सरपंचांची मिरवणूक} पाहायला.

%%K t03-c2-11-1 p
११. --- सर्वांत आधी, अहो, काही \VUgras{मुले} \textsuperscript{(35)}
\VUgras{वादकांच्या} \textsuperscript{(36)} पुढे चालत आहेत; मग \VUgras{सरपंच}
\textsuperscript{(37)}, त्यांचे सहकारी आणि गावाची \VUgras{नगरपालिका सभा} येते.
शेवटी \VUgras{अग्निशामक} \textsuperscript{(38)} चालतात.

%%K t03-tit-12 sub
\VUsaut{5.0mm}
\VUpk{10.2pt}{40}{जत्रा.}
\VUsaut{3.6mm}

%%K t03-c2-12-1 p
१२. --- \VUgras{जत्रेच्या मैदानावर} \textsuperscript{(39)} पुष्कळ माणसे आहेत.
निरनिराळे \VUgras{तंबू} पाहायला लोक येतात: \VUgras{लाकडी घोडे}
\textsuperscript{(40)}, \VUgras{सर्कस} \textsuperscript{(41)} --- जिथे खेळ सुरूही
झाला आहे ---, आणि \VUgras{सार्वजनिक नाच} \textsuperscript{(42)}, जिथे गावातले तरुण
आणि तरुणी \VUgras{नाचण्याचा} आनंद घेतात.

%%K t03-c2-13-1 p
१३. --- पलीकडे आपल्याला \VUgras{आखाडा} \textsuperscript{(44)} दिसतो, जिथे शूर
स्पॅनिश \VUgras{झुंजार}

%%K t03-c2-13-1 p suite
पिसाळलेल्या \VUgras{बैलाशी} \textsuperscript{(45)} झुंजतो; मग \VUgras{घसरगुंडी}
\textsuperscript{(49)}, आणि \VUgras{नेमबाजीची जागा} \textsuperscript{(46)}, जिथे लोक
\VUgras{बंदुका} चालवायचा आणि \VUgras{नेमावर} मारायचा सराव करतात.

%%K t03-c2-14-1 p
१४. --- जवळच, अहो, \VUgras{जत्रेतले नाट्यगृह} \textsuperscript{(47)}, जिथे
\VUgras{विनोदिका} आणि \VUgras{शोकांतिका} सादर केल्या जातात. अगदी जवळ
\VUgras{महाकाय माणसाचा} \textsuperscript{(48)} आणि \VUgras{बुटक्याचा} तंबू आहे.
पहिल्याची \VUgras{उंची} दोन मीटर पंधरा सेंटिमीटर आहे; दुसरा जेमतेम एका मुलाएवढा
आहे.

%%K t03-c2-15-1 p
१५. --- शेवटी, अहो, मोठा \VUgras{प्राणिसंग्रह} \textsuperscript{(50)}, त्याच्या
\VUgras{रानटी जनावरांसह}, त्याच्या बलदंड \VUgras{नखांच्या} \VUgras{वाघासह}
\textsuperscript{(51)}, दाट \VUgras{आयाळीच्या} \VUgras{सिंहासह}
\textsuperscript{(52)}, त्याच्या दोन \VUgras{जिराफांसह} \textsuperscript{(53)} आणि
त्याच्या \VUgras{हत्तीसह} \textsuperscript{(54)}, जो आपली \VUgras{सोंड} इतक्या
कौशल्याने वापरतो.

%%K t03-c2-16-1 p
१६. --- या जनावरांना शिकवणारा \VUgras{माणसाळवणारा} \textsuperscript{(55)}
तंबूच्या दारात आहे.

\VUsaut{6.0mm}
\VUfilet{22mm}
